\chapter{Required prior knowledge}

\section{Assumed prior knowledge}
The target audience are engineering students.
We assume Bachelor level engineering knowledge of mathematics,
physics, programming and control.
The list of required topics is given below,
followed by some key concepts from said topic.
%A simple litmus test for whether you know enough is the following.
%Say I give you a URL to a csv file which contains
%the points of a n-dimensional quadratic function.
%Can you need to find the minimum of this function?
%This would entail that you open an empty Python file, 
%download the file and load the points, 
%do a least-squares fit of a quadratic on them,
%and then either implement vanilla gradient descent 
%or explicitly compute the minimum by solving for x at gradient 
%equal to zero (and knowing why both of these work.)

\begin{itemize}
		\item (multivariate) calculus: partial derivatives, gradients,...
		\item linear algebra: matrix inverse, nullspace...
		\item mechanics: dynamics, Lagrangian mechanics
		\item basic convex optimization: convex functions, 
				gradient descent,...
		\item basic programming: OOP concepts, algorithm complexity,...
		\item introductory control course: error dynamics,  
				stability, PID control,...
\end{itemize}


\section{Topics covered in pre-requisites tutorials}
In order to work with robots, you will need to some
very basic operating systems (OS) concepts.
The target OS is Linux as it the standard in robotics,
but the topics covered apply to any OS.
In order to work with the library, you will need to know
how to program in Python and some experience
with a few important packages such as NumPy.
Since these are not always covered outside of a
computer or data science curriculum,
dedicated tutorials have been developed to fill in the gaps.
All of these topics are covered in the context of
the SMC library and robotic control with concrete examples.

Following the pre-requisites there are dedicated
tutorials on using SMC. Finally, 
there is a small crash introduction to Lie algebra for robotics,
and a crash introduction to project management as
even small robotics projects can get very convoluted,
very fast, if one is not careful and deliberate.

Tutorials on Linux, Python and the SMC library 
are in 
\href{https://gitlab.control.lth.se/marko-g/LPR-tutorials}{LPR tutorials}.
These tutorials are in both text and video formats 
(TODO atm the text is only a list of topics, while
the videos are actual tutorials).
Just the videos can be found in the following 
\href{https://www.youtube.com/playlist?list=PLKOlTxWSIaGWsaJwz6oKtgGB1Ukv_RIeP}{YouTube playlist}.

The tutorials aim to give you an overview of the topics,
tell you only the essential information, and refer you
to other resources for anything else. You might need
all of them, or none of them --- have a look to be sure
you're covered.
The topics covered, and the main goal are:
\begin{itemize}
		\item The GNU/Linux command line interface (CLI) - 
				navigating directories, installing packages
		\item Docker - (special) virtual machines,
				building images (SMC Docker image) and running
				containers
		\item Networking - IP address and ports, creating
				new connections (creating a manual fixed-IP connection
				to the robot), basic debugging
		\item Python - syntax, standard library and some gotchas,
				basic packages like NumPy, matplotlib, SciPy
		\item Multiprocessing - most basic inter-process communication
				(IPC), running a function in parallel in Python
		\item Computer vision - common modern approaches
				and short introduction to OpenCV
		\item Robotics - high level description
				of what a robot is and basic control 
				approaches
		\item Lie algebra - concepts required to understand
				robot kinematics - \textbf{it's easier
				to learn this than to try to avoid it}
		\item SMC - intro and overview, Pinocchio
				(package for robot math),
				tutorial on leveraging exiting controllers,
				tutorial on writing a new controller
		\item Project management - project planning,
				iterative prototyping, software management with git
\end{itemize}
There is also bonus content for some of the topics
to try to convince you that it's worth it to dig
a bit deeper.

Books, websites and documentation of other topics will be provided 
as they appear in the text.


\section{Robotics}
\label{sec-robotics-into}
This section is by absolutely no means a proper
introduction to robotics. 
For this the reader is referred to the following books:
\cite{murray,lynch}.
In particular the following chapters:
\begin{itemize}
		\item TODO: Lynch, list key chapters
		\item TODO: Murray, list key chapters
		\item \href{https://scaron.info/robotics/jacobian-of-a-kinematic-task-and-derivatives-on-manifolds.html}{this excellent blog post}.
\end{itemize}

%The twist $ V_{ \text{error} }  $ 
%corresponds to the velocity such that moving 
%from frame $ \left\{ e \right\}    $, with twist $ V_{  e}  $, for
%one unit of time, leads to the $ \left\{ g \right\}   $ 
%frame.
%Said another way, if $ \left\{ V_{  e} \right\}   $ is
%the body velocity of the end-effector frame for one unit
%of time, it gets to the goal frame.
%The reason this works has to do with Lie algebra.
%Detailed overview can be found in ex. 
%Frank C. Park, Kevin Lynch - Modern Robotics
%or Murray, R.M., Li, Z. and Sastry, S.S., 2017. A mathematical introduction to robotic manipulation. CRC press,
%or on 
%\href{https://scaron.info/robotics/jacobian-of-a-kinematic-task-and-derivatives-on-manifolds.html}{this excellent blog post}.

However, we can not allow perfection to be the enemy of good.
Going through a robotics course and reading a few introductory 
books is objectively speaking not a hard requirement 
to get into robotics and robot control.
This is because we can leverage software to hide and 
automatically deal with many of the robotics concepts
and math. Just as one does not need to have read the
TCP standard and the rest of RFCs in order to open
a web browser and search the internet.
Having that said, working with robots is much more
involved than searching for recipes online,
and some amount of technical detail is simply unavoidable.
This is especially the case for engineers working with robots.
Also, knowledge is power, and the more knowledge you have,
the more you can do.

The purpose of this section is to at least lay out
the basic concepts and terminology required to specify 
a task for the robot. Additional mathematical detail
is laid out in the Control section \ref{sec-control}.
Depending on what you're trying to achieve you'll need
none of that, or it won't be nearly enough.
Hence it is advised to follow a top-down approach, 
i.e. to utilize the concepts described here and the 
ready-made controllers until this stops being enough the problem at
hand, at which point one should dig into the underlying math
to get a good understanding of what's happening.

\subsection{Basic terminology}
\label{sec-robotics-terminology}
A robot is series of \textbf{actuators} (motors) connected with 
\textbf{links}. For the purposes of this text, all
links are \textbf{rigid} and the motors work perfectly
(no friction, instantaneous reference reaching).
A robot \textbf{joint} is a combination of an actuator
and a link --- basically the robot's equivalent of limb.
In the case of a rotational motor we have
a revolute joint whose position is described in radians.
A series of joints creates a \textbf{kinematic chain}.
If only one end is ``bolted'' to a table or floor,
while the other is free, we call it an \textbf{open kinematic chain}.
Otherwise, it is a \textbf{closed kinematic chain}.
An open kinematic chain robot consisting of revolute joints
is referred to as a \textbf{robotic manipulator} (more or less).
Thus one might consider a quadruped a combination of 4 manipulators,
which form a closed kinematic chain when legs are touching the ground.

Since we have access to joint positions through encoders
in the motors, 
the robot is most naturally modelled with Lagrangian mechanics.
The set of all joint positions $ q  $ 
forms \textbf{joint/configuration space}.
The physical 3D environment around the robot
is referred to as \textbf{operational space}.

Robots are controlled either through \textbf{velocity control}
$ u = \dot{q}  $
or \textbf{torque control} $ u = \tau  $. 
If the robot moves relatively slowly, it's acceleration
is low, and thus the dynamic effects are small.
In such a case velocity control suffices and dynamic
effects can simply be treated as errors.
In the majority of the text this is the assumed scenario.
However, if the robot is carrying heavy loads, has to move
quickly, or needs to balance (ex. if it's a humanoid),
then its full dynamics need to be modelled 
and the robot has to be controlled through torque control.

Since we want the robot to act in its environment,
we usually define tasks in operational space.
But we have to control the robot in joint space!
Thus it is essential to understand the relationship between
joint and operational space. How this is done then
directly informs the control strategy. 
Hence, we postpone this discussion to the Control section
\ref{sec-control}.

As it turns out, the operational space is not so simple.
There are two main problems which need addressing.
The first is the definition of a position of an object.
A point in space can be described with 3 numbers corresponding
to x,y, and z coordinates. However, an object is more than just
one point! Thus we need \textbf{position} and \textbf{orientation}
to describe where a (rigid) object is in space.
Critically, orientations (rotations) do not form a vector space,
meaning they can not be added or subtracted directly.
Thankfully this is not a problem at all if one uses the \textbf{Lie
algebra} formalism, which is introduced in the following subsection.
The second problem is that every single relevant quantity
in the environment has to be described in some \textbf{coordinate frame}.
And as it very quickly turns out, trying to directly
define everything in one coordinate frame is cumbersome
to the point of impossibility.
The solution is to define quantities in convenient coordinate
frames, and define \textbf{coordinate transformations} 
between these frames.
Thankfully, this approach is very easy once one gets used 
to it (although will always and forever remain somewhat tedious).

\subsection{Geometry}
\label{sec-robotics-geometry}
\subsubsection{Rigid body coordinates}
A point in space $ p \in \mathbb{R}^{3}  $ can be describe
with 3 coordinates. A rigid body has many points,
but they are always in the same relation to each other,
no matter where the body is in space.
Hence we can define a coordinate frame in, and attached to,
the body itself. If we know where this coordinate frame
is, we know where every point of the body is.
This is thus exactly equivalent to describing
where one coordinate frame is with respect to another.
And, once we know we have this description,
we also have the recipe to transform coordinates of 
these coordinate frames.

The position of the coordinate frame is simply the position
of its origin, so again $  p \in \mathbb{R}^{3}    $.
Its orientation can be determined by describing each unit axis
of the frame, which yields a \textbf{rotation matrix}
$ R \in \mathbb{R}^{3 \times 3}  $
\begin{equation}
		R = \begin{bmatrix} \hat{x} & \hat{y} & \hat{z} \end{bmatrix} .
\end{equation}
Alternatively, one can use quaternions or Euler angles
to represent rotation.
On paper, or in code, the choice doesn't matter as long as the software
transforms whatever representation you gave it into
quaternions, and handles all the subsequent operations for you,
which is the case in Pinocchio (and thus in SMC).
Long story short, 
rotation matrices and quaternions correctly
cover the space of rotations (and are equivalent descriptions), 
and quaternions are more
computationally efficient.
Euler angles do not map all rotations correctly,
but are the only truly intuitive description ---
they allow you to rotate something by 20 degrees.

Here there are 2 key things to note.
Firstly, in practise, one can always \textbf{visualize all the frames and
operations on them} and visually verify that the correct 
operational is taking place.
Secondly, for purposes of designing control,
we only care about the \textbf{mathematical properties}
of rotations, and never care about the exact values.
Thus, the idea is to know the theory, code it up correctly,
and then set the frames correctly once, and verify them visually.
\textbf{With computers, there is absolutely no benefit in flipping
through coordinates in your head.}

With this in mind, we will develop the theory by
assuming the rotation matrix representation.
Here we should note the following.
A position of rigid body is given by its position $ p  $ and orientation $ R  $.
Orientations can be written down in many ways,
and in this text it will be a rotation matrix $ R  $.
The pose (coordinate frame) of the body $  \left\{ b \right\}     $ 
in frame $  \left\{ a \right\}     $ is given by
$ T_{ a }^{ b }  $.
The same notation is used for rotations: $ R_{ a }^{ b }  $
is the orientation of frame $ \left( b \right)   $ as seen in frame $ \left( a \right)   $.
Mathematically,
unless explicitly stated otherwise, this object should be considered as a 
homogeneous  transformation matrix 
\begin{equation}
T_{ a }^{ b }= 
\begin{bmatrix} 
	R_{ a }^{ b } \in \mathbb{R}^{3 \times 3} & p_{ a }^{ b } \in \mathbb{R}^{3} \\
	0 \in \mathbb{R}^{1 \times 3} & 1
\end{bmatrix} 
\end{equation}
The reason for this format is the following.
If you multiply 2 of these matrices, 
the result will again be a homogeneous transformation matrix.
And why do we care about that? Because multiplying
homogeneous transformation matrices is how we do transforms
coordinates from one frame to another!
For example, say you have the position of frame
$ \left\{ c \right\}   $ in frame $ \left\{ b \right\}   $,
i.e. $ T_{ b }^{ c }  $, and likewise you have the position
of frame $ \left\{ b \right\}   $ in frame $ \left\{ a \right\}   $,
i.e. $ T_{ a }^{ b }  $.
To get the position of frame $ \left\{ c \right\}   $ in 
coordinates of frame $ \left\{ a \right\}   $ we simply do
\begin{equation}
T_{ a }^{ c } = T_{ a }^{ b } T_{ b }^{ c }
\end{equation}
And yes, you can always cancel corresponding superscripts and subscripts 
like this!
Lastly, similar to rotations,
\begin{equation}
		(T_{ a }^{ b })^{ -1 } = T_{ b }^{ a }
\end{equation}

Equipped with this knowledge, you are ready to define a target
for the robot to go to, and just apply a controller 
to reach this target.
\textbf{You don't need any extra knowledge to start programming the robot!}

However, to develop new controllers for the robot, and get the most
out of the existing ones, we need to go one step further.

\subsubsection{Lie algebra for robotics}
\label{sec-robotics-lie-algebra}
As mention, the space of rotations is not a vector space,
so we can't add rotations together.
The space of 3D rotations is denoted by $ SO (3)  $.
In mathematics $ SO (n)  $ groups are called orthogonal groups,
and they fall into the class of groups called Lie groups.
There are different equivalent representations of orthogonal groups,
and rotation matrices are one of them.
Rotation matrices have the following properties:
\begin{gather}
		R^{ T } R = I  \implies R^{ -1 } = R^{ T } \\
		det(R) = 1
\end{gather}
And now we can write that
for $ R_{ 1 }, R_{ 2 } \in SO(3), R_{ 1 } + R_{ 2 } \notin SO(3)  $.
Mathematically speaking, $ SO(3)  $ is not a field,
but a group. If using the rotation matrix representation,
the group operation $ (\cdot)  $ becomes matrix multiplication,
and for quaternions it becomes quaternion multiplication.
The point of this is that we can abstractly describe what the group is,
and what the group operation is (every operation has an inverse, there is
an identity element and so on).
Then any representation that complies with the abstract description works.

But what we care about for robotics is the following.
Say you have some $ R_{ 1 }  $ and $ R_{ 2 }  $,
and how what to know what is the ``difference'' between them.
As in how much do you need to rotation $ R_{ 1 }  $ to get to $ R_{ 2 }  $.
This is computed as follows:
\begin{equation}
R_{ 1 }^{ 2 } = R_{ 1 }^{ -1 } R_{ 2 } = R_{ 1 }^{ T } R_{ 2 }
\end{equation}

Rigid body transformations, coordinate transformations,
and coordinate frame descriptions are all just differently
interpreted elements of the special Euclidean group $ \mathbb{SE} (3)  $,
which is simply the combination of the orthogonal group and 3-dimensional
vectors space: $  \mathbb{SE} (3) = \mathbb{R}^{3} \times SO (3)   $.
It is likewise a Lie group.
It follows the same story as rotations:
$    T_{ 1 }, T_{ 2 } \in \mathbb{SE}(3), T_{ 1 } + T_{ 2 } \notin \mathbb{SE}(3)  $
for exactly the same reason.
And we compute the ``difference'' between two frames in the same way:
\begin{equation}
T_{ 1 }^{ 2 } = T_{ 1 }^{ -1 } T_{ 2 } = T_{ 1 }^{ T } T_{ 2 }
\end{equation}

\paragraph{So what are rigid body velocities then?}
Defining a vector field in a vector space is a straightforward affair.
Every point in space $ p \in \mathbb{R}^{n}  $ gets an associated vector 
$ v \in \mathbb{R}^{n}  $ of the same dimension.
The set of these vectors again form a vector space.
In particular, if you have a vector $ v_{ 1 }  $ at point $ p_{ 1 }  $,
you can add or subtract it from a vector $ v_{ 2 }  $ at point $ p_{ 2 }  $
and have a well-defined operation.
To ground this in physics, say there is a particle following some trajectory
$ \gamma (s) \in \mathbb{R}^{ 3 }, s \in \mathbb{R}  $.
Tangents of curve $ \gamma (s)  $, i.e. $ \dot{\gamma} (s)  $ would be these vectors,
and they would represent the instantaneous velocity at each point of the curve.
In this case, what is implied by the vector field being a vector space is that
we can directly compare velocities at different points. A vector 
$ v = \begin{bmatrix} 2 & 1 \end{bmatrix}   $ will mean exactly the same
thing be it attached to point $ p_{ 1 }  $ or $ p_{ 2 }  $.

But Lie groups are not vector spaces. Formally speaking, they form a complete
manifold, which is to say they form a ``curved space''. 
Understanding what this really means requires learning differential geometry.
But here we simply want to have a conceptual understanding to help
us apply correct equations and understand what is going on with a robot.
So in lieu of mathematical formalism, consider the following example.
The surface of the Earth is a sphere. While space is 3-dimensional,
we only need to coordinates to uniquely describe every point on the Earth:
latitude and longitude. Thus the surface of Earth is a 2-dimensional object.
This 2-dimensional object is curved. When you add two of its points described in 3D,
you leave the surface and end up somewhere in space.
Thus viewing it as an object defined in 3D vector space is not helpful.

Now imagine there is an airplane travelling the surface of the Earth along a curve $ \gamma(s)  $,
going ``straight'' from Australia to Great Britain.
Imagine the 3D velocity vector of the airplane when it is above Australia.
And now imagine how it looks when it's above Britain.
They are pointing in different directions. But, from the pilot's perspective,
she has been flying in the same direction, and measured the same velocity
through the whole flight!
Thus, in this scenario, it also isn't helpful to consider velocities 
as 3D vectors in ``ambient space'', as we can not meaningfully compare
velocities at different points of the sphere.


So now what?
The first thing to note is that at each point $ p  $,
the set of velocities does form a vector space, because clearly
it makes sense to compare velocities while being at the same point.
This vector space is called a \textbf{tangent space}, and it is always
associated with the point it is attached to.
The second thing to note is that the velocity vector changes 
smoothly between $ \gamma (s)  $ and $ \gamma (s + \epsilon)  $.
So one can form a differential relationship between 
$ \dot{\gamma} (s)  $ and $ \dot{\gamma} (s + \epsilon)  $.
This of course depends on how the space in which $ \gamma  $ lives is curved,
but the point is a sensible mathematical description can be constructed.

Turning back to the airplane example. The airplane is travelling at a constant velocity,
which is some vector $ v  $. Looking from ambient space, it changes direction
a we progress along $ \gamma (s)  $, while looking the same from the curved surface.
Moving a constant vector along a curve on a curved surface 
is done via \textbf{parallel transport}.
If we know what the surface is, and know how to do parallel transport,
we can know where we will end up, denoted by $ p_{ 2 }  $, if we start at $ p_{ 1 }  $
and go with velocity $ v  $ for a unit period of time.
To cut a long story short, this is done via
the \textbf{exponential map}:
\begin{equation}
p_{ 2 } = \exp (v) p_{ 1 }
\end{equation}
Likewise, we can compute a velocity one needs to travel at to go from
$ p_{ 1 }  $ to $ p_{ 2 }  $ in a unit amount of time via
the \textbf{logarithmic map}:
\begin{equation}
v = \log (p_{ 1 }^{ 2 }) = \log (p_{ 1 }^{ -1 }p_{ 2 })
\end{equation}
Note that for rotation matrices and homogeneous transformation matrices,
the operations $ \exp  $ and $ \log  $ are the matrix exponential
and matrix logarithm respectively (they're different for quaternions),
and there are efficient algorithms for their computation.

Here are some important properties to keep in mind.
To begin, the set of velocities forms a \textbf{Lie algebra}.
Formally speaking, it is defined as the tangent space of the identity
element of the associated Lie group. 
Thus a Lie algebra is a vector space.
Speaking loosely, we can meaningfully compare different tangent vectors (velocities)
by mapping them to the same tangent space via parallel transport,
and then mapping them back to their original points.
Note that this is mathematical construction baked into the exponential
and logarithmic maps, and you don't need to compute this yourself.
For the $ SO (3)  $ group, its Lie algebra is denoted with 
$ so (3) \subset \mathbb{R}^{3} $,
and for the $ \mathbb{SE}(3)  $ group its Lie algebra
is denote with $ se(3) \subset \mathbb{R}^{6}  $.

Crucially for us, the exponential and logarithmic maps provide a direct
unique correspondence between a Lie group and its Lie algebra.
What does this mean? The exponential map turns a velocity,
i.e. a tangent vector into a group element, and the logarithmic map does the opposite:
\begin{gather}
T_{ 2 } = T_{ 1 }^{ 2 }T_{ 1 } \\
\log:  T_{ 1 }^{ 2 }\mapsto v_{ 1 } \\
\exp: v_{ 1 } \mapsto T_{ 1 }^{ 2 } 
\end{gather}

With this, we can return to rigid body kinematics.
Say you have an object at location
$ T_{ w }^{ o }  $ described in the world frame,
and you want to know how it should move to get to $ T_{ w }^{ goal }  $.
First you can compute the transformation from the current pose to the goal pose:
\begin{equation}
T_{ o }^{ goal } = (T_{ w }^{ o })^{ -1 } T_{ w }^{ goal }
\end{equation}
Taking $ v_{ o } = \log (T_{ o }^{ goal })  $, we get that 
if the object, described by its coordinate frame $ \left\{ o \right\}   $,
travels at velocity $ v_{ o }  $ for one unit of time,
it will reach $ T_{ w }^{ goal }  $ from its starting position.
The velocity coordinates corresponding to positions are called
\textbf{linear velocities} $ v  $, while those corresponding to orientations
are called \textbf{rotational velocities} $ \omega  $.
When we want to emphasize that a velocity is an element of $ se (3)  $,
we call it a \textbf{twist} and write $ \mathcal{V} = 
\begin{bmatrix} v & \omega \end{bmatrix}^{ T } \in se (3) $.
\textbf{Note:} it is crucial to note that the velocity $ v_{ o }  $
is \textbf{described}  in the $ \left\{ o \right\}   $ frame in this calculation.
This is maybe a bit weird to think about as the frame $ \left\{ o \right\}   $
is moving together with the object. In other words, while that velocity is constant,
it is changing with respect to a static frame.
While somewhat confusing, this does not pose a problem when defining 
control laws which leverage this description.

\paragraph{What about forces?}
Again, to cut a long story short, forces form a vector space
in a manner very similar to velocities.
The positional part of a force vector is called a \textbf{linear force} $ f  $,
while the rotational part is called a \textbf{torque} $ \tau  $.
When we want to emphasize that we are treating forces with
the Lie group formalism, we refer to them as \textbf{wrenches},
$ \mathcal{F} = \begin{bmatrix} f & \tau \end{bmatrix} ^{ T }   $.
Formally speaking, wrenches are elements of the dual space of 
Lie algebra denote $ se*(3)  $, which is also a vector space.
This is important to have consistent definitions of work and power,
but neither will be encountered in this text.

\subsubsection{Note on practical use}
In SMC, all kinematics and dynamics math is handled by the Pinocchio library.
The code is written with the Lie algebra formalism.
So there is a class \fbox{\texttt{SE3}} which describes an $ \mathbb{SE}(3)  $
object. It can be constructed from a rotation matrix and a position vector,
or a quaternion or a set of Euler angles. Likewise, 
it can output any of these representation back.
Every of the hitherto mention operations is implemented:
group operation $ (\cdot)  $, inverses, exponential and logarithmic maps,
and other operations that this text skipped (such as mapping a velocity to different
frames with the adjoint map).

All of these operations use quaternions in the background,
which are the most computationally efficient representation.
Furthermore, all of this happens in highly optimized C++ code which is exposed
to Python via bindings.
You are \textbf{strongly encouraged} to rely on these operations.
This offers the following benefits:
\begin{itemize}
		\item This ensures the fastest possible computations on the market,
				meaning your controller will be able to meet real-time requirements.
		\item You do not need to think of representation details,
				but only need to focus on the Lie algebra abstractions 
				while implementing the math. This dramatically decreases the
				mental load during implementation.
\end{itemize}

This is the reason why it is not necessary to learn all the mathematics
to be able to program a robot using model-based control techniques, including developing
new controllers --- all the ugly robotics math is hidden away behind an API. 
From an practical operational standpoint, you are \textbf{strongly encouraged}
to invest a few hours of days programming and visualizing different
transformation and motions. What you want to achieve is that the following
3 things correspond to each other harmoniously:
\begin{enumerate}
		\item math on paper
		\item math in code
		\item the visualization of that math in your head
\end{enumerate}
In other words, you want to achieve a visual correspondence of a motion in your head 
and on the screen,
i.e. you put down an equation on paper, then in code, thinking it looks like X,
and then it really looks like X when visualized.
Once you achieve this state you can easily define any task in operational space.

Again, having that said, knowledge is power, and if you want to get serious
about robot control, you will have to read the books, and be able to derive
all the math results from first principles: 
there is no other way to \textit{truly and completely} understand
what is going on.
From the author's personal experience, everything clicked only after
also implementing all the equations in software as well --- programming
exposes all handwaving in a constructive derivation.

%\section{Software}
%There are dependencies so important and integral to the SMC library
%that they should be learned and gotten used to in their own right.
%While they are quite common, some users may encounter them for the first time.
%It is of course impossible to cover everything in detail,
%but we can at least guide you to where you can learn more
%and provide some context to make this a more pleasant experience
%with a clear end in sight.
%The aim of this section is to:
%\begin{itemize}
%		\item briefly introduce SMC's dependencies
%		\item explain their role in SMC
%		\item refer to their respective documentation
%				and pointing to the most relevant parts
%		\item give some non-obvious but helpful tips and tricks
%\end{itemize}
%
%\subsection{Operating systems (Linux)}
%\subsubsection{Basic of the Command Line Interface}
%
%\subsubsection{Docker}
%
%\subsubsection{Parallel computing}
%
%\subsubsection{Networking}
%
%\subsection{Python}
%\subsubsection{NumPy}
%\subsubsection{Matplotlib}
%\subsubsection{OpenCV}
%\subsubsection{Pinocchio}
%\subsection{Optional dependencies}
%\subsubsection{Crocoddyl}
%\subsubsection{Pink}
%
%\section{Project management}
%\subsection{Project plans}
%\subsection{Git}
